\begin{lemma}[The intrinsic anticanonical surface family]\hypertarget{tag-032B}{}\label{push:src:D.1.2}
\StatementRecord{push:src:D.1.2}{proved}{032B}{the definitions and constructions stated in the proof}

Let \(V\) be a connected Riemann surface and
\(\pi:\mathcal X\to V\) a flat proper holomorphic map of reduced
complex spaces. Suppose that over a nonempty open subset \(V_0\), the
fibres are smooth compact threefolds whose unique anticanonical divisor
is \(2S_v\), where \(S_v\) is a smooth connected surface with
\(b_1(S_v)=2\). Then there is an open subset \(V^{\rm good}\subset V\), with at
most countable complement, on which the reduced support of the
relative anticanonical divisor is a smooth proper surface family. Its
Albanese elliptic invariant is holomorphic.
\end{lemma}


\begin{proof}
\begin{enumerate}[label=\arabic*.,leftmargin=*]
\item Flatness and the smoothness of the fibres over \(V_0\) imply that
\(\pi\) is a submersion there: the local analytic smoothness criterion
applies because the base is smooth. The relative nonsmooth locus is a
closed analytic subset, by the Jacobian criterion. Its proper image is analytic by the coherent direct-image argument in Theorem~\ref{night:P2.4}, applied to the structure sheaf of this proper analytic locus.
It misses \(V_0\), hence is discrete. We remove it and work over the
connected remaining curve \(V'\). Removing a discrete subset from a
connected Riemann surface does not disconnect it. Thus
\(\mathcal X'\to V'\) is a smooth proper map.

\item Write \(L=K_{\mathcal X'/V'}^{-1}\). Grauert's direct-image theorem \cite[\S6, Hauptsatz~I, p.~59]{Grauert60} makes \(F=\pi_*L\) and \(R^1\pi_*L\) coherent. His erratum \cite[item~2, p.~36]{Grauert1963ErratumD} corrects the assertion in \S7, Satz~2 that the locus where the comparison map fails to be bijective is always closed. For the present smooth curve, the following exact sequence computes the canonical base-change map and its torsion obstruction.
For a local base parameter \(t\) vanishing at \(v\), smoothness and
invertibility of \(L\) give an exact sequence
\(0\to L\xrightarrow{t}L\to i_*L_v\to0\), where
\(i:X_v\hookrightarrow\mathcal X'\) is the closed embedding.
Its direct-image sequence gives, at \(v\),
\[
 0\longrightarrow F_v/tF_v\longrightarrow H^0(X_v,L_v)
 \longrightarrow (R^1\pi_*L)_v[t]\longrightarrow0.
\]
The last term denotes elements annihilated by \(t\). The support of
the torsion subsheaf of the coherent module \(R^1\pi_*L\) is closed
discrete, by the structure theorem for finitely generated modules
over the local discrete valuation rings of the smooth base curve.
Outside that support the canonical base-change map is an isomorphism.
Also \(F\) is torsion-free: multiplication by a nonzero base function
is injective on sections of \(L\). It is therefore locally free.
Its rank is one by choosing an interior point outside the discrete
torsion support.
Locally a frame gives a relative section \(\sigma\) of \(L\),
nowhere identically zero on a fibre. Its divisor is independent of the
frame. Let \(D\) be its reduced support. On \(V_0\), the divisor is
\(2S_v\) and \(D\) is exactly \(S_v\).

\item The map \(D\to V'\) is proper, since \(D\) is closed in a proper
space. There is no vertical divisor component: a vertical divisor in
the smooth total space would contain a whole irreducible fibre,
contrary to the nonzero restriction of \(\sigma\). Every horizontal
component has proper image all of the connected base. Since the
interior fibre is irreducible, there is exactly one horizontal
component, and its multiplicity in \(\operatorname{div}\sigma\) is
two. Thus \(\operatorname{div}\sigma=2D\) on this base. In the smooth
ambient space \(D\) is Cartier. It is flat over the curve because
its local rings are torsion-free over the base discrete valuation
ring. Restriction of the Cartier identity to an interior fibre gives
\(2D_v=2S_v\), whence \(D_v=S_v\) with its reduced structure.
Flatness and the smoothness of these fibres make \(D\to V'\) smooth
over \(V_0\). Its analytic nonsmooth locus therefore misses \(V_0\);
the proper mapping theorem once more makes its image discrete.
After deleting that image, \(D\to V^{\rm good}\) is a smooth proper
family of connected surfaces. Each deletion is closed discrete in the
then-current open curve. Their finite union is at most countable,
and the final good locus is open.

\item Ehresmann's theorem identifies the smooth fibres along every path;
in particular their first Betti number is two. \cite[Theorem~11, p.~297]{Buchdahl1999} then gives a K\"ahler metric on every fibre. Hence
\(h^{1,0}=1\), every holomorphic one-form is closed, and the integral
period map has a rank-two lattice image. Constant \(h^{1,0}\),
properness, and analytic base change supply a local relative
holomorphic one-form. Locally on a simply connected
base disc, a relative one-form and a basis of the free part of
\(H_1\) give two holomorphic periods. Their ratio lies in one
half-plane after orienting the basis. To see holomorphicity directly,
choose local primitives of the closed relative form on simply
connected fibre charts; their differences are holomorphic functions
of the base, and the finite sum around a transported cycle is its
period. The modular invariant of the ratio is independent of all
choices, so these local invariants glue to a holomorphic function on
\(V^{\rm good}\).
\end{enumerate}
\end{proof}
